\section{Avvio dell'Applicazione}

\subsection{Step 1: Avvia i Container}

Dalla directory radice del progetto, esegui:

\begin{lstlisting}
docker-compose up -d
\end{lstlisting}

Il flag \texttt{-d} esegue i container in modalità detached (background).

Docker procederà a:
\begin{enumerate}
    \item Costruire l'immagine dell'applicazione
    \item Scaricare le immagini di Neo4j e Ollama
    \item Avviare i container
    \item Eseguire i health check
\end{enumerate}

\textbf{Tempo di completamento}: 30-60 secondi (più lungo al primo avvio).

\subsection{Step 2: Verifica che i Servizi siano Attivi}

Attendi 30-60 secondi, poi esegui:

\begin{lstlisting}
docker-compose ps
\end{lstlisting}

Dovresti vedere output simile a:

\begin{lstlisting}
NAME                 STATUS              PORTS
mirofish-offline     Up (healthy)        ...
mirofish-neo4j       Up (healthy)        ...
mirofish-ollama      Up                  ...
\end{lstlisting}

Tutti i servizi devono avere status \textbf{Up}.

\subsection{Step 3: Accedi all'Applicazione}

Una volta che i servizi sono attivi, accedi ai seguenti URL:

\begin{itemize}
    \item \textbf{Frontend}: \href{http://localhost:3000}{http://localhost:3000}
    \item \textbf{Backend API}: \href{http://localhost:5001}{http://localhost:5001}
    \item \textbf{Neo4j Browser}: \href{http://localhost:7474}{http://localhost:7474}
    \begin{itemize}
        \item Username: \texttt{neo4j}
        \item Password: \texttt{mirofish}
    \end{itemize}
\end{itemize}
